% =========================================================
%  Paper: CuscoNodes - Sistema de IA Agencia para Alertas
%         Turisticas en Tiempo Real
%  Plantilla RATE — Revista Argentina de Trabajos Estudiantiles
%  Formato dos columnas, A4, estilo IEEE
%  Compatible con pdfLaTeX (TeXpage)
% =========================================================
\documentclass[10pt,twocolumn,a4paper]{article}

\usepackage[T1]{fontenc}
\usepackage[utf8]{inputenc}
\usepackage[spanish,provide=*]{babel}
\usepackage{times}
\usepackage{geometry}
\usepackage{graphicx}
\usepackage{array}
\usepackage{caption}
\usepackage{amsmath}
\usepackage{amssymb}
\usepackage{lettrine}
\usepackage{microtype}
\usepackage{setspace}
\usepackage{titlesec}
\usepackage{fancyhdr}
\usepackage{lastpage}
\usepackage{footnote}
\usepackage{dblfloatfix}
\usepackage{url}
\usepackage[hidelinks]{hyperref}
\usepackage{listings}
\usepackage{xcolor}
\usepackage{booktabs}

\geometry{
  a4paper,
  top=19mm,
  bottom=30mm,
  left=13mm,
  right=13mm,
  columnsep=4mm
}

\lstset{
  language=Python,
  basicstyle=\ttfamily\fontsize{7}{8}\selectfont,
  keywordstyle=\bfseries,
  commentstyle=\itshape,
  stringstyle=\ttfamily,
  breaklines=true,
  frame=single,
  numbers=left,
  numberstyle=\tiny,
  numbersep=4pt,
  showstringspaces=false,
  tabsize=2,
  captionpos=b
}

\setlength{\parindent}{3.5mm}
\setlength{\parskip}{0pt}

\pagestyle{fancy}
\fancyhf{}
\fancyhead[R]{\thepage}
\renewcommand{\headrulewidth}{0pt}

\titleformat{\section}
  {\normalsize\bfseries\centering\scshape}
  {\Roman{section}.}{0.5em}{}
\titlespacing{\section}{0pt}{6pt}{4pt}

\titleformat{\subsection}
  {\normalsize\itshape}
  {\Alph{subsection}.}{0.5em}{}
\titlespacing{\subsection}{0pt}{4pt}{2pt}

\captionsetup[figure]{
  labelsep=period,
  font={small},
  justification=centering,
  position=below
}
\captionsetup[table]{
  labelsep=newline,
  font={small,sc},
  justification=centering,
  position=above,
  name=TABLA
}

\makeatletter
\renewcommand{\maketitle}{%
  \twocolumn[{%
    \begin{center}
      {\fontsize{24}{28}\selectfont\@title\par}
      \vspace{4mm}
      {\normalsize\@author\par}
    \end{center}
    \vspace{4mm}
  }]%
  \footnotetext{$^*$ Trabajo presentado como actividad academica en la
    asignatura de Inteligencia Artificial, Universidad Andina del
    Cusco.}%
}
\makeatother

\title{CuscoNodes: Sistema de Inteligencia Artificial Agencia\\
       para Alertas Turisticas Multilingue en Tiempo Real}

\author{%
  Honorio Morales Ttito$^*$  Sebastian Zavaleta Luna$^*$  Sayo Michael Torres Caceres$^*$ \hfill Escuela Profesional de Ingenieria de Sistemas \\
  \textit{Universidad Andina del Cusco, Cusco, Peru} \\
  \textit{Patrocinador: Israel Rondan (PDS Viajes)}
}

\date{}

% =========================================================
\begin{document}
% =========================================================

\maketitle
\thispagestyle{fancy}

\noindent\textbf{\textit{Resumen}}---CuscoNodes es un sistema de
inteligencia artificial agente que procesa, traduce y notifica
alertas turisticas en la region de Cusco, Peru. El sistema implementa
un pipeline de tres etapas siguiendo el patron arquitectonico
Supervisor-Worker: (1) Percepcion (scraping de fuentes locales),
(2) Razonamiento (analisis y traduccion multilingue con Gemini API), y
(3) Accion (notificacion multicanal via SMTP y Twilio WhatsApp).
Se desarrollaron seis sprints incrementales que cubren desde la
extraccion de datos hasta el despliegue en produccion con panel de
administracion, autenticacion, metricas en tiempo real y
programacion automatica. El pipeline alcanza una latencia promedio de
5.5 segundos con traducciones concurrentes en ingles, frances y
portugues, preservando toponimias sagradas (Machu Picchu,
Ollantaytambo, Sacsayhuaman, Pisac, Chinchero, Urubamba). El sistema
se despliega en Render con Python 3.12.6 y Gunicorn, y el codigo
fuente esta disponible en un repositorio publico de GitHub.

\smallskip
\noindent\textbf{\textit{Abstract}}---CuscoNodes is an agentic
artificial intelligence system that processes, translates, and
notifies tourist alerts in the Cusco region of Peru. The system
implements a three-stage pipeline following the Supervisor-Worker
architectural pattern: (1) Perception (scraping local sources),
(2) Reasoning (multilingual analysis and translation with Gemini API),
and (3) Action (multichannel notification via SMTP and Twilio
WhatsApp). Six incremental sprints were developed, covering from data
extraction to production deployment with an administration dashboard,
authentication, real-time metrics, and automatic scheduling. The
pipeline achieves an average latency of 5.5 seconds with concurrent
translations in English, French, and Portuguese, preserving sacred
toponyms. The system is deployed on Render with Python 3.12.6 and
Gunicorn, and the source code is available in a public GitHub
repository.

% ----------------------------------------------------------
\section{Introduccion}
% ----------------------------------------------------------

\lettrine[lines=2]{E}{}l turismo en la region de Cusco, Peru,
enfrenta desafios de informacion en tiempo real debido a la
frecuencia de eventos locales como huelgas, cierres de vias y
condiciones climaticas adversas~[1]. Los turistas extranjeros,
principales visitantes de Machu Picchu y el Valle Sagrado,
requieren alertas en su idioma nativo para tomar decisiones
informadas durante su estadia.

CuscoNodes nace como respuesta a esta necesidad, proponiendo un
sistema de inteligencia artificial agente que automatiza el ciclo
completo de deteccion, analisis y comunicacion de alertas. El
proyecto fue desarrollado por el grupo HABA bajo el patrocinio de
PDS Viajes, con el objetivo de integrar tecnologias de IA, APIs
comerciales y canales de mensajeria en un pipeline cohesivo.

El proyecto se encuentra disponible en los siguientes enlaces:
\begin{center}
\small
\url{https://github.com/Honorio-Morales/cusconodes} \\
(GitHub -- codigo fuente)
\end{center}
\begin{center}
\small
\url{https://cusconodes.onrender.com} \\
(Render -- aplicacion desplegada)
\end{center}
\begin{center}
\small
\url{https://honorio-morales.github.io/cusconodes/} \\
(GitHub Pages -- poster academico)
\end{center}

% ----------------------------------------------------------
\section{Fundamento Teorico}
% ----------------------------------------------------------

\subsection{Sistemas Multi-Agente}

Los sistemas multi-agente (MAS) constituyen un paradigma de la
inteligencia artificial distribuida donde multiples agentes
autonomos interactuan para resolver problemas complejos~[2]. Cada
agente posee capacidad de percepcion, razonamiento y accion,
operando de manera independiente pero coordinada.

El patron Supervisor-Worker, utilizado en este proyecto, establece
una jerarquia donde un agente supervisor coordina y delega tareas a
agentes trabajadores especializados, sincronizando sus resultados
antes de continuar~[3]. Este patron es particularmente efectivo
para pipelines de procesamiento donde diferentes etapas requieren
experticia especifica.

\subsection{Agentes de Lenguaje Grande (LLMs)}

Los modelos de lenguaje de gran escala (LLMs) como Gemini~[4] y
GPT~[5] han demostrado capacidad para tareas de traduccion
contextual, clasificacion y generacion de texto. En este proyecto,
Gemini API se utiliza para la traduccion multilingue de alertas,
preservando terminos geograficos sagrados mediante tecnicas de
prompt engineering con system prompts especificos.

\subsection{Arquitectura Supervisor-Worker}

La arquitectura implementada consta de tres roles:
\begin{itemize}
  \item \textbf{SupervisorAgent:} orquestador central que lee
  alertas del directorio de persistencia, evalua la urgencia y
  delega tareas a los workers.
  \item \textbf{TranslatorWorker:} agente especializado en
  traduccion que invoca Gemini API para generar contenido en un
  idioma destino (en, fr, pt).
  \item \textbf{NotifierWorker:} agente de difusion que despacha
  notificaciones via SMTP (Gmail) y Twilio WhatsApp.
\end{itemize}

% ----------------------------------------------------------
\section{Metodologia}
% ----------------------------------------------------------

\subsection{Arquitectura del Sistema}

La solucion se organiza en tres capas principales:
\begin{itemize}
  \item \textbf{Backend (Flask):} servidor REST que expone
  endpoints para orquestacion, consultas, historial, metricas y
  configuracion.
  \item \textbf{Agentes (Python):} modulo de logica agente
  implementando el pipeline Supervisor-Worker con concurrencia
  via \texttt{threading}.
  \item \textbf{Frontend (Tailwind CSS):} dashboard corporativo
  modo claro con visualizacion en tiempo real del pipeline,
  historial, metricas y controles de administracion.
\end{itemize}

La Figura~1 muestra la arquitectura de despliegue del sistema y la
Figura~2 presenta una captura del dashboard operativo.

\begin{figure}[t]
\centering
\framebox[\columnwidth]{\parbox{0.9\columnwidth}{
\small\ttfamily
Cliente Web $\rightarrow$ \\
\hspace*{2mm}Flask/Gunicorn $\rightarrow$ \\
\hspace*{4mm}SupervisorAgent $\rightarrow$ \\
\hspace*{6mm}TranslatorWorker [Gemini API] \\
\hspace*{6mm}NotifierWorker \\
\hspace*{8mm}SMTP (Gmail) \\
\hspace*{8mm}Twilio WhatsApp \\
\hspace*{4mm}data/processed/ (JSON) $\leftarrow$ Persistencia
}}
\caption{Arquitectura de despliegue del sistema CuscoNodes.}
\label{fig:arquitectura}
\end{figure}

\begin{figure}[t]
\centering
\framebox[\columnwidth]{\parbox{0.9\columnwidth}{\centering\vspace*{3cm}
{\small [Captura: Dashboard de CuscoNodes mostrando el pipeline
de agentes (Supervisor, Traductor, Notificador) con estado
FINALIZADO y metricas de desempeno.]}\vspace*{3cm}}}
\caption{Dashboard de orquestacion en tiempo real del pipeline
Supervisor-Worker.}
\label{fig:dashboard}
\end{figure}

\subsection{Tecnologias Utilizadas}

La Tabla~I resume las tecnologias principales del proyecto.

\begin{table}[t]
\caption{Tecnologias del Sistema}
\label{tab:tecnologias}
\centering
\small
\begin{tabular}{p{2.5cm}|p{4cm}}
\hline
\textbf{Componente} & \textbf{Tecnologia} \\
\hline
Backend & Flask, Gunicorn, Python 3.12.6 \\
\hline
Frontend & Tailwind CSS v4, Chart.js v4 \\
\hline
IA & Google Gemini API (gemini-2.0-flash) \\
\hline
Mensajeria & Twilio WhatsApp Business API \\
\hline
Correo & SMTP Gmail con App Password \\
\hline
Scheduler & APScheduler 3.x \\
\hline
Despliegue & Render (free tier) \\
\hline
Control versiones & Git + GitHub \\
\hline
\end{tabular}
\end{table}

\subsection{Pipeline de Agentes}

El pipeline ejecuta cuatro etapas secuenciales:

\textbf{Etapa 1 -- SupervisorAgent:} Lee el archivo JSON mas
reciente del directorio \texttt{data/processed/}. Si la alerta es
\texttt{IRRELEVANTE}, aborta el pipeline.

\textbf{Etapa 2 -- TranslatorWorker $\times$ 3 (concurrentes):}
Tres workers se ejecutan en hilos paralelos, cada uno invocando
Gemini API para traducir a un idioma. Incluyen retry con backoff
exponencial (1s, 2s, 4s) ante errores de cuota (HTTP 429). Si
fallan todos los reintentos, usan traducciones de respaldo.

\textbf{Etapa 3 -- NotifierWorker:} Despacha el payload
multilingue via SMTP (plantilla HTML con banderas y toponimias
destacadas) y Twilio WhatsApp (mensaje Markdown).

\textbf{Etapa 4 -- Persistencia:} Guarda el resultado completo en
\texttt{data/processed/alertas\_clasificadas\_demo\_final.json}.

\subsection{Concurrencia y Sincronizacion}

La implementacion usa \texttt{threading.Thread} para lanzar los
tres workers de traduccion simultaneamente. Se aplica un stagger
de 0.5 segundos entre workers para evitar rate limiting en Gemini
API. La sincronizacion se realiza mediante \texttt{thread.join()},
garantizando que la consolidacion ocurra solo despues de que todos
los workers completen su ejecucion.

\subsection{Proteccion de Toponimias}

El sistema mantiene un conjunto inmutable de terminos protegidos:
Machu Picchu, Ollantaytambo, Sacsayhuaman, Pisac, Chinchero y
Urubamba. Estos se preservan mediante system prompts en las
solicitudes a Gemini API y se resaltan visualmente en el dashboard
y las plantillas HTML de correo.

\subsection{Endpoints REST}

La Tabla~II enumera los endpoints implementados.

\begin{table}[t]
\caption{Endpoints REST del Sistema}
\label{tab:endpoints}
\centering
\small
\begin{tabular}{p{0.8cm}|p{2.6cm}|p{2.2cm}|p{0.7cm}}
\hline
\textbf{Met.} & \textbf{Endpoint} & \textbf{Funcion} & \textbf{Auth} \\
\hline
GET & /api/health & Health check & No \\
POST & /api/login & Autenticacion & No \\
GET & /api/history & Historial (50) & Si \\
GET & /api/metrics & Metricas agregadas & Si \\
GET & /api/export/csv & Exportar CSV & Si \\
POST & /api/orchestrate & Pipeline agentes & Si \\
POST & /api/query & Consulta IA & No \\
GET/POST & /api/recipients & Destinatarios & Si \\
POST & /api/scheduler/* & Control scheduler & Si \\
GET & /api/monitoring & Monitoreo & Si \\
\hline
\end{tabular}
\end{table}

\subsection{Plan de Monitoreo}

El sistema incluye un subsistema de monitoreo con:
\begin{itemize}
  \item \textbf{Health check:} endpoint \texttt{GET /api/health}
  que verifica la disponibilidad del servicio.
  \item \textbf{Logging persistente:} registro de cada ejecucion
  del pipeline en \texttt{data/pipeline.log}.
  \item \textbf{Tracking de ejecuciones:} archivo
  \texttt{data/monitoring.json} con conteo de exitos, fallos,
  timestamps de ultima ejecucion y ultimo error.
  \item \textbf{Dashboard de monitoreo:} panel visual en el
  frontend mostrando estado del scheduler, health status,
  pipeline exitosos/fallidos y ultimo error registrado.
\end{itemize}

% ----------------------------------------------------------
\section{Resultados}
% ----------------------------------------------------------

\subsection{Sprints Completados}

El proyecto se desarrollo en seis sprints incrementales. La
Tabla~III resume el avance por sprint.

\begin{table}[t]
\caption{Avance por Sprint}
\label{tab:sprints}
\centering
\small
\begin{tabular}{p{1.2cm}|p{4.5cm}|c}
\hline
\textbf{Sprint} & \textbf{Enfoque} & \textbf{Estado} \\
\hline
Sprint 1 & Iniciacion y planificacion & 100\% \\
Sprint 2 & Percepcion (scraping RPP, PeruRail) & 100\% \\
Sprint 3 & Razonamiento (agentes mock + dashboard) & 100\% \\
Sprint 4 & Accion (Gemini, SMTP, Twilio reales) & 100\% \\
Sprint 5 & Administracion (auth, metricas, scheduler) & 100\% \\
Sprint 6 & Cierre y documentacion final & 100\% \\
\hline
\end{tabular}
\end{table}

\subsection{Metricas de Rendimiento}

La Tabla~IV presenta las metricas obtenidas del pipeline en
ejecuciones reales.

\begin{table}[t]
\caption{Metricas de Rendimiento del Pipeline}
\label{tab:metricas}
\centering
\small
\begin{tabular}{p{3.5cm}|c}
\hline
\textbf{Metrica} & \textbf{Valor} \\
\hline
Latencia de traduccion paralela (3 idiomas) & 1.5 s \\
Latencia total del pipeline & 5.50 s \\
Toponimias preservadas & 6/6 \\
Hilos de traduccion concurrentes & 3 \\
Canales de notificacion & 2 (WA + SMTP) \\
Codigo HTTP endpoint POST & 200 OK \\
\hline
\end{tabular}
\end{table}

\subsection{Resultados por Historia de Usuario}

La Tabla~V detalla el cumplimiento de las 21 historias de usuario
distribuidas en los sprints 2 a 5. Cada historia fue verificada
mediante criterios de aceptacion documentados en los reportes
tecnicos de cada sprint.

\begin{table*}[t]
\caption{Historias de Usuario Completadas}
\label{tab:historias}
\centering
\small
\begin{tabular}{p{1cm}p{3.8cm}p{5.5cm}c}
\toprule
\textbf{CUN} & \textbf{Historia} & \textbf{Implementacion} & \textbf{Estado} \\
\midrule
01 & Scraping RPP & \texttt{RPPScraper} con parsing de HTML & 100\% \\
02 & Scraping PeruRail & \texttt{PeruRailScraper} con extraccion de tablas & 100\% \\
03 & Base de datos JSON & Persistencia en \texttt{data/raw/} y \texttt{data/processed/} & 100\% \\
04 & Configuracion entorno & \texttt{Settings} con python-dotenv & 100\% \\
05 & Clasificador urgencia & Clasificacion CRITICA/ALTA/NORMAL/IRRELEVANTE & 100\% \\
06 & Filtro relevancia & Filtro turistico con Criticality Guard & 100\% \\
07 & Orquestacion central & \texttt{SupervisorAgent} con pipeline secuencial & 100\% \\
08 & Traduccion concurrente & 3 \texttt{TranslatorWorker} en hilos paralelos & 100\% \\
09 & Notificacion multicanal & \texttt{NotifierWorker} con SMTP + WhatsApp & 100\% \\
10 & Dashboard supervisor & Interfaz web con Tailwind CSS modo claro & 100\% \\
11 & Endpoint REST orquestacion & \texttt{POST /api/orchestrate} & 100\% \\
12 & Historial persistente & \texttt{GET /api/history} con JSON persistente & 100\% \\
13 & Panel metricas & Latencia y estado en dashboard & 100\% \\
14 & Toponimias preservadas & System prompt + \texttt{protected\_terms} & 100\% \\
15 & Traduccion Gemini real & \texttt{google.genai} con retry backoff & 100\% \\
16 & Agentes cooperativos & Supervisor + Workers con \texttt{threading} & 100\% \\
17 & SMTP real & Gmail SMTP con plantilla HTML profesional & 100\% \\
18 & Consulta IA & \texttt{POST /api/query} con Gemini & 100\% \\
19 & Twilio WhatsApp real & Sandbox con mensajes Markdown & 100\% \\
20 & Plantillas alerta HTML & Banderas y toponimias destacadas & 100\% \\
21 & Autenticacion y monitoreo & Login, metricas, scheduler, export CSV & 100\% \\
\bottomrule
\end{tabular}
\end{table*}

\subsection{Resultados de Despliegue}

El sistema se desplego exitosamente en Render (free tier) con:
\begin{itemize}
  \item Python 3.12.6 fijado via \texttt{.python-version}
  \item Servidor Gunicorn con 2 workers
  \item Variables de entorno configuradas en dashboard de Render
  \item Health check endpoint \texttt{/api/health} retorna 200 OK
   \item Tiempo de respuesta inicial ~30s (wake-up de free tier)
   \item Credenciales por defecto: usuario \texttt{admin}, contrasena \texttt{cusconodes2025} (configurable via \texttt{ADMIN\_PASSWORD})
   \item Sesiones persistentes con \texttt{FLASK\_SECRET\_KEY} fijo para consistencia entre workers de Gunicorn
\end{itemize}

La Figura~3 muestra el panel de administracion con metricas
en tiempo real, scheduler automatico y monitoreo del sistema.

\begin{figure*}[t]
\centering
\framebox[\textwidth]{\parbox{0.95\textwidth}{\centering\vspace*{4cm}
{\small [Captura: Panel de administracion con metricas en tiempo
real (ejecuciones, tasa de exito, latencia promedio), grafico de
alertas por dia, controles del scheduler automatico y panel de
monitoreo del sistema.]}\vspace*{4cm}}}
\caption{Panel de administracion con metricas, scheduler y
monitoreo del sistema CuscoNodes.}
\label{fig:admin}
\end{figure*}

% ----------------------------------------------------------
\section{Discussion}
% ----------------------------------------------------------

\lettrine[lines=2]{L}{}os resultados confirman la viabilidad de
implementar un sistema de IA agente funcional utilizando APIs
comerciales y un presupuesto de infraestructura minimo (Render free
tier). La arquitectura Supervisor-Worker demostro ser efectiva para
orquestar tareas de traduccion y notificacion de manera
concurrente.

\subsection{Lecciones Aprendidas}

\begin{itemize}
  \item \textbf{Limitaciones de cuota gratuita:} Gemini API impone
  limites de tasa (429) que requieren reintentos con backoff
  exponencial y un sistema de fallback robusto.
  \item \textbf{Sandbox de Twilio:} solo numeros registrados
  manualmente pueden recibir mensajes, lo que limita la
  escalabilidad inmediata.
  \item \textbf{Concurrencia Python:} el \texttt{Global
  Interpreter Lock (GIL)} limita el paralelismo real, pero
  \texttt{threading} es suficiente para operaciones de I/O como
  llamadas API.
  \item \textbf{Render free tier:} el servicio duerme tras 15
  minutos de inactividad, requiriendo ~30s para reactivarse.
\end{itemize}

\subsection{Trabajo Futuro}

Como lineas de trabajo futuro se identifican:
\begin{itemize}
  \item Migracion a un plan de pago de Gemini API para eliminar
  los fallbacks por cuota.
  \item Integracion de Twilio con numeros verificados (fuera del
  Sandbox) para envio a destinatarios no registrados.
  \item Implementacion de una base de datos SQL (PostgreSQL) para
  reemplazar la persistencia en JSON plano.
  \item Autenticacion multi-usuario con roles (admin, lector).
  \item Notificaciones push via Firebase Cloud Messaging.
\end{itemize}

% ----------------------------------------------------------
\section{Conclusiones}
% ----------------------------------------------------------

\begin{enumerate}
  \item Se implemento un sistema completo de IA agente que
  automatiza el ciclo de deteccion, traduccion y notificacion de
  alertas turisticas en Cusco.
  \item El pipeline Supervisor-Worker con traduccion concurrente
  en tres idiomas alcanza una latencia promedio de 5.5 segundos,
  cumpliendo el requisito de menos de 8 segundos.
  \item Las 21 historias de usuario (CUN-01 a CUN-21) fueron
  implementadas y validadas, cubriendo percepcion, razonamiento,
  accion y administracion.
  \item La integracion de APIs reales (Gemini, SMTP, Twilio)
  reemplazo los mocks del Sprint 3, demostrando la viabilidad
  tecnica del sistema en un entorno de produccion.
  \item El despliegue en Render con Python 3.12.6 y Gunicorn
  proporciona una plataforma accesible para demostraciones y
  validacion con el patrocinador.
  \item El sistema preserva la integridad de las toponimias
  sagradas mediante system prompts y validacion visual en el
  frontend.
\end{enumerate}

% ----------------------------------------------------------
\section*{Enlaces del Proyecto}
% ----------------------------------------------------------

\begin{itemize}
  \item \textbf{Codigo fuente:}
  \url{https://github.com/Honorio-Morales/cusconodes}
  \item \textbf{Aplicacion en vivo:}
  \url{https://cusconodes.onrender.com}
  \item \textbf{Poster academico:}
  \url{https://honorio-morales.github.io/cusconodes/}
\end{itemize}

% ----------------------------------------------------------
\section*{Comandos de Reproduccion}
% ----------------------------------------------------------

\begin{lstlisting}[language=bash, numbers=none, frame=single,
  basicstyle=\ttfamily\fontsize{8}{9}\selectfont]
# Clonar repositorio
git clone https://github.com/Honorio-Morales/cusconodes.git
cd cusconodes

# Crear entorno virtual
python3 -m venv venv && source venv/bin/activate

# Instalar dependencias
pip install -r requirements.txt

# Copiar y editar variables de entorno
cp .env.example .env
# Credenciales por defecto: admin / cusconodes2025

# Iniciar servidor local
python web_server.py
# Abrir http://127.0.0.1:5000
\end{lstlisting}

% ----------------------------------------------------------
\section*{Reconocimientos}
% ----------------------------------------------------------

Se agradece a la Escuela Profesional de Ingenieria de Sistemas de
la Universidad Andina del Cusco por el espacio academico para el
desarrollo de este trabajo, y al patrocinador Israel Rondan de PDS
Viajes por la orientacion en las necesidades del dominio turistico.

% ----------------------------------------------------------
\section*{Referencias}
% ----------------------------------------------------------
\begingroup
\small
\begin{enumerate}
\setlength{\itemsep}{0pt}
\setlength{\parsep}{0pt}

\item Ministerio de Comercio Exterior y Turismo del Peru,
  ``Perfil del Turista Extranjero 2024,'' MINCETUR, Lima, Peru,
  2024.

\item M. Wooldridge, \textit{An Introduction to MultiAgent
  Systems}, 2nd ed. Chichester, UK: Wiley, 2009.

\item E. H. Durfee, ``Distributed problem solving and planning,''
  in \textit{Multi-Agent Systems}, G. Weiss, Ed. Cambridge, MA:
  MIT Press, 2013, ch. 3, pp. 69--130.

\item Google DeepMind, ``Gemini: A family of highly capable
  multimodal models,'' 2024. [Online]. Available:
  \url{https://arxiv.org/abs/2312.11805}

\item A. Vaswani et al., ``Attention is all you need,''
  in \textit{Proc. NeurIPS}, pp. 5998--6008, 2017.

\item J. Dean and S. Ghemawat, ``MapReduce: Simplified data
  processing on large clusters,'' in \textit{Proc. OSDI},
  pp. 137--150, 2004.

\item G. E. Box and D. R. Cox, ``An analysis of transformations,''
  \textit{Journal of the Royal Statistical Society, Series B},
  vol. 26, no. 2, pp. 211--252, 1964.

\item Python Software Foundation, ``threading -- Thread-based
  parallelism,'' Python 3.12 Documentation, 2024. [Online].
  Available: \url{https://docs.python.org/3/library/threading.html}

\item Render Inc., ``Python version,'' Render Docs, 2026.
  [Online]. Available: \url{https://render.com/docs/python-version}

\end{enumerate}
\endgroup

\end{document}
